% !TeX root = ../数值分析与科学计算笔记.tex
% 由主文档通过 \input 载入；请编译主文档。
\setcounter{chapter}{4} % 保留原书章号。
\chapter{最小二乘、正交分解与函数逼近}
\label{chap:na-ls}

\section{超定方程与正交投影}
测量数据往往给出比未知数更多的方程，噪声使它们无法同时满足。给定 $A\in\mathbb R^{m\times n}$、$b\in\mathbb R^m$ 且 $m\ge n$，\emph{线性最小二乘（linear least squares）}求
\begin{equation}
 \min_{x\in\mathbb R^n}\|Ax-b\|_2^2.
 \label{eq:na-ls-objective}
\end{equation}
这里优化变量为 $x$，数据为 $A,b$，目标是残差平方和，而非插值所有数据。

\begin{theorem}[正规方程与唯一性]
\label{thm:na-normal-equations}
向量 $x_*$ 是式~\eqref{eq:na-ls-objective} 的极小点，当且仅当
\begin{equation}
 A^{\mathsf T}Ax_*=A^{\mathsf T}b.
 \label{eq:na-normal-equations}
\end{equation}
式~\eqref{eq:na-normal-equations} 称为正规方程（normal equations）。最小二乘解总存在；若 $A$ 列满秩则唯一，此时 $A^{\mathsf T}A$ 对称正定。
\end{theorem}
\begin{proof}
最优点处沿任意方向 $z$ 的一阶导数为 $2z^{\mathsf T}A^{\mathsf T}(Ax_*-b)=0$，故正规方程必要。反之，若 $r_*=b-Ax_*$ 满足 $A^{\mathsf T}r_*=0$，则
\[
 \|A(x_*+z)-b\|_2^2=\|r_*\|_2^2+\|Az\|_2^2,
\]
故它为极小点。有限维列空间是闭子空间，对 $b$ 的正交投影存在，投影向量可表示为 $Ax_*$，从而解存在。列满秩时非零 $z$ 满足 $Az\ne0$，故极小点唯一且正规矩阵正定。
\end{proof}

正规方程是理论刻画，不代表数值上总应显式形成它。若 $A$ 列满秩，则 $\kappa_2(A^{\mathsf T}A)=\kappa_2(A)^2$，因此形成正规方程会放大条件方面的困难。这里矩形矩阵的二范数条件数按最大、最小奇异值之比定义。

\begin{example}[直线拟合]
对数据 $(0,1),(1,2),(2,2)$，拟合 $p(t)=\alpha+\beta t$。设计矩阵与正规方程为
\[
 A=\begin{pmatrix}1&0\\1&1\\1&2\end{pmatrix},\quad
 A^{\mathsf T}A=\begin{pmatrix}3&3\\3&5\end{pmatrix},\quad
 A^{\mathsf T}b=\begin{pmatrix}5\\6\end{pmatrix}.
\]
解得 $\alpha=7/6$、$\beta=1/2$。按 $b-Ax$ 约定，残差为 $(-1/6,1/3,-1/6)^{\mathsf T}$，与设计矩阵两列都正交，残差平方和为 $1/6$。拟合曲线没有穿过每一个点，这是目标函数允许的行为。
\end{example}

\section{QR 分解与 Householder 变换}
列满秩矩阵可写为约化 QR 分解 $A=Q_1R$，其中 $Q_1\in\mathbb R^{m\times n}$、$Q_1^{\mathsf T}Q_1=I$，$R\in\mathbb R^{n\times n}$ 为非奇异上三角矩阵。将 $Q_1$ 扩充为正交矩阵 $Q=(Q_1,Q_2)$，利用正交变换保二范数，有
\[
 \|Ax-b\|_2^2=\|Rx-Q_1^{\mathsf T}b\|_2^2+\|Q_2^{\mathsf T}b\|_2^2.
\]
所以只需求解 $Rx=Q_1^{\mathsf T}b$。当 $m=n$ 时第二项为空项，视为零。

\emph{Householder 反射（Householder reflection）}定义为 $H=I-2vv^{\mathsf T}/(v^{\mathsf T}v)$，其中 $v\ne0$。直接相乘可验证 $H^{\mathsf T}=H$、$H^2=I$。给定非零向量 $z$，取 $\alpha=-\operatorname{sign}(z_1)\|z\|_2$，约定 $\operatorname{sign}(0)=1$，再取 $v=z-\alpha e_1$，即可使 $Hz=\alpha e_1$。符号选择避免 $v$ 首分量由两个接近的大数相减得到。

逐列把反射作用于剩余子矩阵即可得到 QR 分解，无须显式存储每个完整 $H$；更新可写为矩阵与向量的运算。稠密 $m\times n$ 矩阵且 $m\ge n$ 的代价约为 $2mn^2-2n^3/3$。Householder QR 在标准浮点条件下有范数意义的后向稳定性。经典 Gram--Schmidt 正交化在接近线性相关时可能丢失正交性，修正版本和重新正交化能改善这一现象；这些算法得到相同精确分解的思想，却有不同有限精度表现。

\section{奇异值分解与正则化}
\begin{theorem}[奇异值分解的基本形式]
任意 $A\in\mathbb R^{m\times n}$ 可写为 $A=U\Sigma V^{\mathsf T}$，其中 $U,V$ 为正交矩阵，$\Sigma$ 为矩形对角矩阵，非负对角元 $\sigma_1\ge\cdots\ge\sigma_r>0$，其余为零，$r=\operatorname{rank}(A)$。这些正数称为奇异值（singular values）。
\end{theorem}
\begin{proof}[证明要点]
对称半正定矩阵 $A^{\mathsf T}A$ 有正交特征向量组 $v_i$，非负特征值可写为 $\sigma_i^2$。对 $\sigma_i>0$ 定义 $u_i=Av_i/\sigma_i$，可验证这些向量正交归一。把它们补成 $\mathbb R^m$ 的正交基，并对零特征值对应的核空间选取正交基，即得分解。这里使用线性代数中的实对称谱定理。
\end{proof}

由奇异向量展开，最小二乘解中二范数最小者为
\[
 x^\dagger=\sum_{i=1}^r\frac{u_i^{\mathsf T}b}{\sigma_i}v_i.
\]
在 $\ker A$ 方向增加任意向量不会改变残差，却增大该最小范数解的范数。很小的 $\sigma_i$ 会放大数据中对应方向的扰动，形成逆问题中的不稳定性。

\emph{Tikhonov 正则化（Tikhonov regularization）}在本书采用参数平方约定：给定 $\lambda>0$，求
\[
 \min_x\bigl(\|Ax-b\|_2^2+\lambda^2\|x\|_2^2\bigr),\qquad
 x_\lambda=\sum_{i=1}^r\frac{\sigma_i}{\sigma_i^2+\lambda^2}(u_i^{\mathsf T}b)v_i.
\]
新的目标有唯一解，因为 $A^{\mathsf T}A+\lambda^2I$ 正定。正则化通过减弱小奇异值方向引入偏差并降低扰动放大；选择 $\lambda$ 需要噪声水平、验证数据或其他先验信息，不能仅以训练残差最小来决定。

\section{连续最小二乘与正交多项式}
设 $w$ 在 $[a,b]$ 上连续且严格为正，定义 $\langle f,g\rangle_w=\int_a^b f(x)g(x)w(x)\,\mathrm dx$。给定连续函数 $f$ 与有限维子空间 $V=\operatorname{span}\{\phi_0,\ldots,\phi_n\}$，最佳平方逼近求 $p\in V$ 使 $\|f-p\|_w^2=\langle f-p,f-p\rangle_w$ 最小。若基函数线性无关，则 Gram 矩阵 $G_{ij}=\langle\phi_j,\phi_i\rangle_w$ 正定，存在唯一解，且 $f-p$ 与 $V$ 正交。证明与定理~\ref{thm:na-normal-equations} 相同，只需把欧氏内积替换为积分内积。

当基函数两两正交时，系数直接为 $c_j=\langle f,\phi_j\rangle_w/\langle\phi_j,\phi_j\rangle_w$。在 $[-1,1]$、$w=1$ 下，前几个 Legendre 多项式是 $1,x,(3x^2-1)/2$。选择正交基改善系数求解的表示条件，但连续最佳平方逼近仍不同于一致范数下的最佳逼近，也不同于有限样本拟合。

\begin{example}[平方函数的最佳一次平方逼近]
在 $[-1,1]$ 上以权 $1$ 逼近 $f(x)=x^2$。由于奇偶性，$x$ 的系数为零；常数项为 $\int_{-1}^1x^2\,\mathrm dx/\int_{-1}^1 1\,\mathrm dx=1/3$。故 $p(x)=1/3$，平方误差为 $\int_{-1}^1(x^2-1/3)^2\,\mathrm dx=8/45$。
\end{example}

\section{习题}
\begin{enumerate}[label=\arabic*.]
\item\label{exr:na-ls-constant} 用常数拟合 $b_1,\ldots,b_m$，证明最小二乘解为样本平均数；给出正权 $w_i$ 下的加权结果。
\item\label{exr:na-ls-householder} 证明 $H=I-2vv^{\mathsf T}/(v^{\mathsf T}v)$ 保持二范数，且将 $v$ 映为 $-v$。
\item\label{exr:na-ls-rank} 求 $A=\begin{pmatrix}1&1\\2&2\end{pmatrix}$、$b=(1,2)^{\mathsf T}$ 的全部最小二乘解及最小范数解。
\item\label{exr:na-ls-projection} 求 $x^3$ 在 $[-1,1]$、权 $1$ 下的最佳一次平方逼近及平方误差。
\end{enumerate}

